\clearpage
\subsection{匹配圈}\label{sec:18h}
现在更仔细地考察上述情形的一个重要特例，即 $L=\Sigma_{\mu,I}$ 是 $S$ 中某条匹配路径 $\mu$ 所对应的匹配圈，而 $\mu$ 相对于各 $\gamma_k$ 的位置如\sr{fig:18.1}{图~18.1} 所示。给 $\widetilde\Delta_k,\widetilde\Delta_{k+1}$ 赋予等变膜结构，使 $V_k,V_{k+1}$ 上诱导的膜结构彼此同痕；这只是对\stref{引理}{18.7}中固定分次方式的重述。类似地，选择 $L$ 的一个分次，并给 $\widetilde L$ 赋予相应的等变膜结构。

\StatementNumber{18}\begin{lemma}\label{st:18.18}
$HF^*(\widetilde\Delta_k^\#,\widetilde\Delta_{k+1}^\#)\cong H^*(S^n;\K)$，其上的 $\Z/2$ 作用平凡。
\end{lemma}
\begin{proof}
这只是\eref{18.23}的特例：$HF^*(\widetilde\Delta_k^\#,\widetilde\Delta_{k+1}^\#)\cong HF^*(V_k^\#,V_{k+1}^\#)$，且 $\Z/2$ 作用平凡；后一群同构于 $H^*(V_k;\K)$，因为两个膜彼此同痕。
\end{proof}

\StatementNumber{19}\begin{lemma}\label{st:18.19}
群 $HF^*(\widetilde L_\pm^\#,\widetilde\Delta_k^\#)$ 都是一维的，且集中于同一个次数。
\end{lemma}
\begin{proof}
仔细考察匹配圈的构造，就应当能看清这一点。利用\stref{引理}{15.3}，可找到一个沿各纤维恒定的变形，从给定的纤维化结构出发，最终到达一个可用常值同痕 $I$ 按朴素方式定义匹配圈的 Lefschetz 纤维化。可要求该变形集中在 $\mu$ 本身的原像附近，因此它提升为 $\widetilde E$ 的同类变形。在变形后的结构中，所考察的 Lagrangian 子流形的对应物，记为 $\widetilde L_\pm^*$ 与 $\widetilde\Delta_k^*$，横截相交于唯一一点；该点是 $E_{\gamma_k(1)}$ 中临界点在 $\widetilde E$ 中的一个原像。由\stref{引理}{7.1}，变形后的二重分歧覆盖带有到原来 $\widetilde E$ 的共形
\SourcePage{293}
辛嵌入；在同痕意义下，它将 $\widetilde L_\pm^*$ 映到 $\widetilde L_\pm$，将 $\widetilde\Delta_k^*$ 映到 $\widetilde\Delta_k$。第一种情形显然成立，因为这类同痕本来就是匹配圈构造的一部分；第二种情形则可由考察整个变形过程中连续产生的消失圈族看出。因此，可对所考察的 Lagrangian 子流形作同痕，使它们横截相交于唯一一点，这当然说明 Floer 上同调是一维的。为说明非平凡部分位于同一次数，只需使用 $\Z/2$ 作用 $\eps$；按定义，它交换这两个 Floer 群。
\end{proof}

为简便，假设 $L$ 的分次已选成使群 $HF^*(\widetilde L_\pm^\#,\widetilde\Delta_k^\#)$ 集中于次数 $1$。

\StatementNumber{20}\begin{lemma}\label{st:18.20}
在 $H(\widetilde{\A})$ 中，有如下形式的正合三角：
\SourceEquationDiagram{diagram-18-26}{0.238845897}{18.26}
\end{lemma}
\begin{proof}
由\stref{引理}{18.7}及\stref{推论}{17.17}，在 $D(\widetilde{\A})$ 中有
\begin{equation}\label{eq:18.27}
\widetilde\Delta_{k+1}^\#\cong T_{\widetilde L_-^\#}T_{\widetilde L_+^\#}(\widetilde\Delta_k^\#).
\end{equation}
此外，$\widetilde L_\pm$ 不相交，故两者之间没有态射，从而
\[
T_{\widetilde L_-^\#}T_{\widetilde L_+^\#}(\widetilde\Delta_k^\#)
\cong T_{\widetilde L_-^\#,\widetilde L_+^\#}(\widetilde\Delta_k^\#),
\]
其中右端的定义见\stref{笔记}{5.1}、\stref{笔记}{5.2}。于是，\eref{18.27}导出 $H(\Tw\widetilde{\A})$ 中如下形式的正合三角：
\SourceDiagram{diagram-18-h-triangle}{0.616845490}
由\stref{引理}{18.19}，左上顶点是 $S^{-1}(\widetilde L_-^\#)\oplus S^{-1}(\widetilde L_+^\#)\cong S^{-1}(\widetilde L^\#)$。现在旋转这个三角即可，见\stref{推论}{3.17}。\znote{285}{$S^{-1}$ 是 §12a 的平移函子（分次减 $1$，并张量 $\lambda^{\mathrm{top}}$ 的逆），它出现是因为引理 18.19 把两个群都定在次数 $1$，而三角的左顶点必须写成对象的直和。图中两条边上的 $[1]$ 同样是分次平移，不是同调次数加一：它记录态射 $\widetilde L^\#\to\widetilde\Delta_k^\#$ 的次数。另注意第三项取 $\widetilde\Delta_{k+1}^\#$ 而非其平移，旋转三角（推论 3.17）之后，方向才与引理所述形式一致；这类「平移指数是否要跟着改」的记账，是本节最易出错之处。[来源：原书 §12a、引理 18.19、推论 3.17（印刷页 285）]}
\end{proof}

\StatementNumber{21}\begin{proposition}\label{st:18.21}
在 $D\F(\pi)$ 中，$\widetilde L^\#$ 同构于 $\Cone(\widetilde\Delta_k^\#\to\widetilde\Delta_{k+1}^\#)$，其中锥可在任意非平凡的零次态射上构造。
\end{proposition}
\begin{proof}
\SourcePage{294}
由\eref{18.26}，在 $D(\widetilde{\A})$ 中得到同构 $\widetilde L^\#\cong\Cone(\widetilde\Delta_k^\#\to\widetilde\Delta_{k+1}^\#)$。在同构意义下，将任意态射乘以 $\K^*$ 中的某个元素不会改变其锥。在这里，态射群是一维的，因此只需证明所用的态射非零。若它为零，两个直和项 $\widetilde L_\pm^\#$ 就必分别同构于 $\widetilde\Delta_k^\#,\widetilde\Delta_{k+1}^\#$，但这立即与前面的 Floer 上同调计算矛盾，例如\stref{引理}{18.19}。最后，由于 $\widetilde L^\#$ 是自由的，可应用\stref{引理}{6.2}，得到 $D(\widetilde{\A}^{\mathrm{inv}})=D\F(\pi)$ 中的同构。
\end{proof}

\StatementNumber{22}\begin{note}\label{st:18.22}
另一种方法是将此情形严格视为\stref{命题}{18.17}的特例。为简便，设 $k=1$。此时需考察适当的局部模型，证明乘积
\[
HF^*(\widetilde\Delta_2^\#,\widetilde L_+^\#)\otimes HF^n(\widetilde\Delta_1^\#,\widetilde\Delta_2^\#)
\longrightarrow HF^{*+n}(\widetilde\Delta_1^\#,\widetilde L_+^\#)
\]
是同构；所涉及的各空间实际都是一维的。由此，再结合 $j>2$ 时 $HF^*(\widetilde\Delta_j^\#,\widetilde L_+^\#)=0$，可知\stref{引理}{6.3}证明中的对象 $U$，在这一特殊情形中归约为 $\Cone(\widetilde\Delta_1^\#\to\widetilde\Delta_2^\#)$。
\end{note}
